\documentclass[11pt]{article}
\usepackage{fontspec,amsmath,amssymb,amsthm,hyperref}
\usepackage[margin=1in]{geometry}
\setlength{\emergencystretch}{3em}
\title{The complete canonical KLMC theorem on arbitrary finite coordinate sets}
\author{LangevinLean mathematical exposition}
\date{}
\begin{document}\maketitle
\noindent Source module: \texttt{KineticLangevin/GenericCanonicalKLMC.lean}. Compile this proof document with XeLaTeX or LuaLaTeX. The accompanying HTML pairs the mathematical argument with the exact Lean statements and records the scope of the formalization.


The entire actual-source KLMC proof is stated for a general finite coordinate type \(I\). This includes the particle index \(\{\theta\}\sqcup(\{1,\ldots,N\}\times\{x\})\) without a separate coordinate-isometry transport. All initial couplings, canonical Brownian inputs, selected reference flows, Gibbs stationary laws and genuine Wasserstein infima are the same concrete constructions as in the finite-dimension presentation; the dimension in the final bias is exactly \(|I|\).


\section{Specify the concrete probability space and source projection}


The phase variable is velocity first. The sample space contains the two coupled initial positions, one shared Gaussian velocity and one vector Brownian source. The measure is the actual initial-coupling product extension; probability follows from its first marginal and from the Gaussian/Brownian probability factors. The Brownian-source projection preserves its measure. Pullback along that projection gives a continuous linear map on L\ensuremath{{}^{2}}, and its representative theorem identifies it almost everywhere with evaluation on the actual product sample.



\noindent Lean declarations: \texttt{H}, \texttt{Sample}, \texttt{probability}, \texttt{probability\_isProbabilityMeasure}, \texttt{couplingProbability}, \texttt{brownianLift}, \texttt{brownianLift\_coe}.

\section{Construct the common Brownian input and filtered innovation}


The scaled source lift includes the physical amplitude \(\sqrt{2\gamma}\). At iteration \(n\), the input is the actual vector Brownian increment on \((nh,nh+t]\), and the filtered input is the exponential response to that same increment. The unscaled innovation pair consists of the exact velocity and position Wiener responses at the step endpoint. Both input paths are continuous in L\ensuremath{{}^{2}}, and the additive input starts at zero. These follow from the source lift's continuity and linearity and the proved vector-response properties.



\noindent Lean declarations: \texttt{inputLift}, \texttt{input}, \texttt{filtered}, \texttt{innovation}, \texttt{continuous\_input}, \texttt{continuous\_filtered}, \texttt{input\_zero}.

\section{Identify the exact noise equations and implemented numerical step}


The constructed filtered input satisfies \(B_n(t)=W_n(t)-\gamma\int_0^tB_n(s)ds\). Its endpoint is the scaled velocity innovation, and its actual time primitive is the scaled positional innovation. Lean transports the vector Brownian equations through the continuous linear source lift and commutes that lift with Bochner integration. Induction on the actual recursions then proves that the exponential-input chain is exactly \texttt{klmcChain}, with the canonical correlated innovation pair and the prescribed fluctuation-dissipation scaling.



\noindent Lean declarations: \texttt{filtered\_integral\_equation}, \texttt{integral\_filtered}, \texttt{filtered\_endpoint}, \texttt{exponentialInputChain\_eq\_klmcChain}.

\section{Connect Brownian L\ensuremath{{}^{2}} innovations with the generic output-law construction}


The L\ensuremath{{}^{2}} innovation pair has the same representatives as the pointwise canonical Brownian innovation. The proof pulls back the previously established vector-response representative theorem to the coupling product and combines the two component identities. Comparing with the generic product-law layer's representatives gives equality of the actual L\ensuremath{{}^{2}} pairs. Thus the numerical chain analyzed here has the coupling-independent output law already proved for the implemented KLMC chain.



\noindent Lean declarations: \texttt{innovation\_coe}, \texttt{innovation\_eq\_extensionInnovation}.

\section{Build the actual numerical and stationary-reference initial states}


The two initial phases share exactly one Gaussian velocity variable. The numerical phase uses the first coupled position; the reference phase uses the second. The reference representative is proved to be the explicit pair \((V,X_\mathrm{ref})\). Product independence and the two marginal laws prove that its actual law is \(\mathcal N(0,I_{|I|})\otimes\pi\). The finite target second moment and finite coupling cost supply the required L\ensuremath{{}^{2}} initial positions.



\noindent Lean declarations: \texttt{initialVelocity}, \texttt{numericalInitial}, \texttt{referenceInitial}, \texttt{referenceInitial\_coe}, \texttt{referenceInitial\_law}.

\section{Construct all three chains and segment families}


The numerical chain applies the actual lifted original gradient and the constructed correlated innovations. The reference chain stitches exact selected kinetic segments, each driven by the actual canonical Brownian increment for that iteration. The corresponding reference segments start from these actual reference states. The auxiliary exact segment starts instead from the actual numerical state and receives the identical additive input. The drift is the genuine pointwise gradient lifted to L\ensuremath{{}^{2}}; no substitute potential on the L\ensuremath{{}^{2}} space is assumed.



\noindent Lean declarations: \texttt{numericalChain}, \texttt{referenceChain}, \texttt{referenceSegment}, \texttt{auxiliarySegment}.

\section{Prove every analytic premise of the KLMC coupling theorem}


The bundled specification supplies continuity of auxiliary paths, reference paths and filtered inputs; exact initial states; reference endpoint matching; the implemented exponential update; both actual common-input integral equations; the exponential mild identity; the kinematic primitive identity; and equal initial velocities. Each field follows from the concrete definitions and earlier proved equations. The mild and kinematic fields are derived internally. The only remaining probabilistic input for a final stationary-error theorem is the stationary law of the actual global reference diffusion, proved separately and connected to these stitched segments by uniqueness.



\noindent Lean declarations: \texttt{segments\_spec}.

\section{Construct the full canonical input and identify its increments}


The full input is the scalar canonical Brownian value assembled spatially and pulled onto the actual initial-coupling extension with amplitude \(\sqrt{2\gamma}\). It is continuous in L\ensuremath{{}^{2}} because each assembly and pullback is continuous linear. The Brownian value at zero vanishes: Lean verifies that its deterministic time indicator has empty support, then uses linearity of the isonormal map. For nonnegative local time, the constructed segment input equals the full input difference at \(nh+t\) and \(nh\). The proof checks the nonnegative real-to-NNReal casts and uses linearity of spatial assembly and source pullback.



\noindent Lean declarations: \texttt{globalInput}, \texttt{continuous\_globalInput}, \texttt{globalInput\_zero}, \texttt{input\_eq\_globalInput\_increment}.

\section{Construct the actual full reference solution}


The global reference is the selected kinetic solution with the genuine lifted gradient, the canonical reference initial phase and the full canonical Brownian input. Its specification proves continuity, its exact starting phase, and the actual kinetic integral equation over the whole finite horizon. There is no stationary-law requirement in its definition.



\noindent Lean declarations: \texttt{globalReferenceFlow}, \texttt{globalReferenceFlow\_spec}.

\section{Identify every stitched reference point by uniqueness}


The generic stitching theorem now applies because the local inputs are proved to be actual increments of the full input. Consequently \(R_n=Y_{nh}\), and \(Z_n(t)=Y_{nh+t}\) throughout each step whose endpoint lies in the global horizon. Lean instantiates the existing exact uniqueness results with the actual L\ensuremath{{}^{2}} gradient and checks the horizon arithmetic. Thus a stationary law for the full constructed kinetic solution supplies the law of all reference states used by the quantitative KLMC proof.



\noindent Lean declarations: \texttt{referenceChain\_eq\_globalReferenceFlow}, \texttt{referenceSegment\_eq\_globalReferenceFlow}.

\section{Project the coupling source onto an independent phase/Brownian product}


The projection sends the full sample to \(((V,X_\mathrm{ref}),\omega_\mathrm{Brownian})\). Its law is \((\mathcal N(0,I_{|I|})\otimes\pi)\otimes P_\mathrm{Brownian}\). Lean first projects the initial coupling onto its second marginal, retaining the independent Gaussian/Brownian source. It then applies measure-preserving product associativity and swaps the position and velocity factors. Thus independence of the complete reference initial phase and the complete Brownian source is proved for the actual product construction.



\noindent Lean declarations: \texttt{initialPhase\_brownian\_jointLaw}.

\section{Identify the actual full-input OU increment}


The deterministic integration-by-parts input used in the kinetic splitting proof is \(\Delta W-\gamma\int_0^he^{-\gamma(h-s)}\Delta W_sds\). The actual local-input increment identity replaces each difference of the full canonical input by its constructed segment input. Applying the vector exponential-response formula through the continuous linear source lift proves that this OU increment is exactly the filtered velocity response. Its representative is consequently the scaled canonical vector response on the original Brownian source, pulled onto the coupling extension.



\noindent Lean declarations: \texttt{ouInputIncrement\_globalInput}, \texttt{ouInputIncrement\_globalInput\_coe}.

\section{Derive the full independent thermal Gaussian joint law}


The vector OU stream already has the exact iid Gaussian law with scale \(\sqrt{1-\psi_0(h)^2}\). The product-source joint-law theorem combines this stream with the initial phase target. Composition with the proved phase/Brownian projection law transfers that product law to the actual coupling extension. Finally, the representative identities identify its initial phase with the constructed reference initial L\ensuremath{{}^{2}} phase and its stream with the actual full-input OU increments. The resulting joint law is precisely the source obligation of the proved kinetic Gibbs stationarity transfer theorem.



\noindent Lean declarations: \texttt{initial\_ouInput\_jointLaw}.

\section{Fix the actual numerical output law}


The definition uses the original pointwise KLMC method, the standard Gaussian initial velocity and the exact canonical correlated Brownian innovations. The numerical L\ensuremath{{}^{2}} chain is identified with the generic product-law construction by equality of the actual innovation variables. The established output-law theorem then proves that its pushforward is independent of which initial position coupling was chosen.



\noindent Lean declarations: \texttt{klmcLaw}, \texttt{numericalChain\_outputLaw}.

\section{Transfer one actual global stationary phase law to the reference segments}


The named stationary property concerns the full constructed kinetic flow: at every time its phase law is \(\mathcal N(0,I_{|I|})\otimes\pi\). It contains no numerical convergence conclusion. Exact chain-to-global identification transfers its position marginal to every reference grid point. Exact segment-to-global identification transfers its Gaussian velocity marginal to every intermediate time in a reference segment. Lean obtains these marginals by composition with the product-measure projections.



\noindent Lean declarations: \texttt{GlobalReferenceStationary}, \texttt{referenceChain\_position\_law}, \texttt{referenceSegment\_velocity\_law}.

\section{Apply the proved coupling estimate to one actual initial plan}


The initial position L\ensuremath{{}^{2}} discrepancy equals the true cost of the supplied transport coupling; this follows from the concrete extension's representative and norm identities. The segment specification furnishes all analytic local-error premises internally. Together with the derived stationary position and Gaussian velocity marginals, the quantitative L\ensuremath{{}^{2}} KLMC theorem gives \(W_2(\nu_n,\pi)\le\sqrt2(1-3mh/(4\gamma))^n\operatorname{cost}(c)+Mh\sqrt{2|I|}/m\). Output-law equality and initial-cost equality rewrite this into a bound for the actual numerical law.



\noindent Lean declarations: \texttt{initial\_position\_cost}, \texttt{klmc\_bound\_of\_globalStationarity}.

\section{Prove output finiteness and take the transport infimum}


Finite initial Wasserstein distance supplies an actual finite-cost initial coupling. Its numerical and reference states are genuine L\ensuremath{{}^{2}} variables, so their distance gives a finite transport bound on the output law. This proves output finiteness before real-valued transport arithmetic is used. The scalar recurrence estimate proves the contraction factor is nonnegative under the paper's step restrictions. Lean then applies the proved infimum theorem over all finite-cost couplings, obtaining the same estimate with \(W_2(\mu,\pi)\) in place of the cost of a chosen plan. No existence of an optimal coupling is assumed.



\noindent Lean declarations: \texttt{klmcLaw\_wasserstein2\_ne\_top}, \texttt{klmc\_wasserstein\_bound\_of\_globalStationarity}.

\section{Prove Gibbs stationarity of the actual global coupling reference}


The reference begins with standard Gaussian velocity independent of the target Gibbs position. Its full Brownian forcing and its entire filtered innovation stream are actual constructed objects on the initial-coupling extension. Their exact initial/noise product law is already proved. The kinetic stationarity transfer theorem therefore applies directly to the actual integral solution. For a time \(t\) inside a larger horizon, Lean restricts the actual integral equation to \([0,t]\), supplies the continuous input and its zero starting value, and applies the proved splitting-limit stationarity result. This yields the full phase law \(\mathcal N(0,I_{|I|})\otimes\pi\) at every reference time, without a remaining stationarity or Gaussian-marginal hypothesis.



\noindent Lean declarations: \texttt{globalReferenceFlow\_stationary}.

\section{Conclude the source-paper KLMC Wasserstein estimate}


Let \(f\) be differentiable, \(m\)-strongly convex with \(m>0\), and have an \(M\)-Lipschitz gradient. If \(\gamma^2\ge m+M\), \(\gamma>0\), and \(0<h\le m/(4\gamma M)\), the actual KLMC position law satisfies \[W_2(\nu_n,\pi)\le\sqrt2\left(1-\frac{3mh}{4\gamma}\right)^nW_2(\mu,\pi)+\frac{Mh\sqrt{2|I|}}m.\] The initial position distribution may be arbitrary with finite initial Wasserstein distance; its velocity is the explicit independent standard Gaussian. Lean first derives the finite Gibbs second moment from strong convexity. It then supplies the just-proved global-reference stationarity for every initial transport coupling to the complete canonical coupling/infimum theorem. The conclusion has no conditional local-error, reference-process, Gaussian-law, stationary-law or optimal-coupling premise. The smoothness assumptions are expressed by the actual gradient and strong convexity, so the theorem also covers the paper's Hessian-bounded setting through the existing deterministic geometry results.



\noindent Lean declarations: \texttt{klmc\_wasserstein\_bound}.

\end{document}
